
\paragraph*{Integration and Volume Computation. }

In this chapter, we develop efficient algorithms for volume computation
and integration based on a technique called \emph{Simulated Annealing,
}which is also provably efficient for convex optimization. The volume
problem was defined in Chapter (\ref{chap:Elimination-and-Reduction})
(Problem (\ref{problem:volume})). Integration is the following generalization.
\begin{problem}[Integration]
 Given an evaluation oracle for an integrable logconcave function
$f:\R^{n}\rightarrow\R$, a point $x_{0}$ s.t. $f(x_{0})\ge\beta,$
numbers $r,R>0$ s.t. (i) $\E(\norm{x-x_{0}}^{2})\le R^{2}$ (ii)
the level set of $\pi_{f}$ of measure $1/8$ contains a ball of radius
$r$, and $\varepsilon>0$, output a number $A$ s.t. 
\[
(1-\varepsilon)\int f\le A\le(1+\varepsilon)\int f.
\]

As in all problems in oracle models discussed in this book, the complexity
of an algorithm to solve the above problem is given by the number
of oracles calls and the number of arithmetic operations. 
\end{problem}

Table \ref{tab:volume} below summarizes progress on the volume problem
over the past three decades. Besides improving the complexity of volume
computation, each step has typically resulted in new techniques. For
more details, we refer the reader to surveys on the topic \cite{Simonovits03,VemSurvey}.

\begin{table}[h]
\centering{}%
\begin{tabular}{|c|l|c|}
\hline 
Year/Authors & New ingredients & Steps\tabularnewline
\hline 
\hline 
1989/Dyer-Frieze-Kannan \cite{DyerFK89} & Everything & $n^{23}$\tabularnewline
\hline 
1990/Lov�sz-Simonovits \cite{LS90} & Better isoperimetry & $n^{16}$\tabularnewline
\hline 
1990/Lov�sz \cite{L90} & Ball walk & $n^{10}$\tabularnewline
\hline 
1991/Applegate-Kannan \cite{applegate1991sampling} & Logconcave sampling & $n^{10}$\tabularnewline
\hline 
1990/Dyer-Frieze \cite{DyerF90} & Better error analysis & $n^{8}$\tabularnewline
\hline 
1993/Lov�sz-Simonovits \cite{LS93} & Localization lemma & $n^{7}$\tabularnewline
\hline 
1997/Kannan-Lov�sz-Simonovits \cite{kannan1997random} & Speedy walk, isotropy & $n^{5}$\tabularnewline
\hline 
2003/Lov�sz-Vempala \cite{LV2003} & Annealing, hit-and-run & $n^{4}$\tabularnewline
\hline 
2015/Cousins-Vempala \cite{CV2015} (well-rounded) & Gaussian Cooling & $n^{3}$\tabularnewline
\hline 
2017/Lee-Vempala (polytopes) \cite{lee2017convergence} & Hamiltonian Walk & $mn^{\frac{2}{3}}$\tabularnewline
\hline 
2021/Jia-Laddha-Lee-Vempala \cite{JLLV26} & Faster rounding & $n^{3.5}$\tabularnewline
\hline 
\end{tabular}\caption{\label{tab:volume} The complexity of volume estimation, each step
uses $\widetilde{O}(n)$ bit of randomness. The last algorithm needs
$\widetilde{O}\left(mn^{\omega-1}\right)$ steps per iteration while
the rest need $O(n^{2})$ per oracle query.}
\end{table}
We begin with a description and analysis of the DFK algorithm to estimate
the volume of a convex body in high dimension.

\section{Volume Computation: the DFK Algorithm}
\begin{thm}
There exists a randomized algorithm that given any $\varepsilon,\delta>0$,
and any convex body $K$ with probability at least $1-\delta$ solves
the Volume Estimation problem for $K$ using $\mbox{poly}(n,1/\varepsilon,\ln(1/\delta),\ln R)$
oracle queries and arithmetic operations. 
\end{thm}

The next exercise is the general observation that any algorithm that
succeeds with failure probability less than $1/2$ can be converted
to one with failure probability at most $\delta$. 
\begin{xca}
Show that a randomized algorithm that outputs an acceptable estimate
of a real number with failure probability at most $1/4$ can be converted
to one that outputs an acceptable estimate with failure probability
at most $\delta$ by repeating the algorithm independently $O(\ln(1/\delta))$
times and taking the median of the outputs of the independent trials.

Going forward, we ignore the parameter $\delta$, and focus on algorithms
with constant failure probability. We will now show an algorithm for
volume computation (as a first attempt) whose analysis yields a suboptimal
bound on the sample complexity.

The polynomial-time algorithm of Dyer, Frieze and Kannan \cite{DyerFK89}
is based on the following idea. We consider a sequence of convex bodies
$K_{0}=B_{n}\subset K_{1}\subseteq\ldots\subseteq K_{m}=K$, starting
with the ball contained inside the input body $K$ and ending with
$K$ itself. Then the following telescoping formula suggests an algorithm:
\[
\vol(K)=\vol(K_{0})\cdot\frac{\vol(K_{1})}{\vol(K_{0})}\cdot\ldots\cdot\frac{\vol(K_{m})}{\vol(K_{m-1})}.
\]

To estimate the volume of $K$, we estimate each of the ratios on
the RHS and output their product together with the volume of the unit
ball. We need to set up the sequence so that the sequence is not too
long and this estimator is accurate and efficient. In the DFK algorithm,
each body in the sequence is a ball intersected with the given convex
body $K$:
\[
K_{i}=2^{\frac{i}{n}}B\cap K.
\]
 The length the sequence is $m=\lceil n\log R\rceil$ so that the
final body is just $K$. We can now state the algorithm. In each iteration
the algorithm estimates the ratio $\vol(K_{i})/\vol(K_{i+1})$ by
sampling uniformly from $K_{i+1}$ and using the fraction of the sample
that lies in $K_{i}$ as the estimator.
\end{xca}

\begin{algorithm2e}[H]

\caption{$\mathtt{DFK\ Volume}$}

\SetAlgoLined
\begin{enumerate}
\item For $i=0,\ldots,m=\lceil n\log_{2}R\rceil$, let $K_{i}=2^{\frac{i}{n}}B\cap K.$
\item For $i=0,\ldots,m-1$, 
\begin{enumerate}
\item uniformly sample $X^{(1)},X^{(2)},\ldots,X^{(N)}$ from $K_{i+1}$.
\item Set $Y_{i}=|{\{j:X^{(j)}\in K_{i}\}}|/N$. 
\end{enumerate}
\item Let $Y=\prod_{i=1}^{m-1}Y_{i}$ and return $\vol(B_{n})/Y$.
\end{enumerate}
\end{algorithm2e}

The following observation is immediate.
\begin{lem}
Let $r_{i}=\vol(K_{i})/\vol(K_{i+1})$. Then $\E(Y_{i})=r_{i}$ and
$\Var(Y_{i})=r_{i}(1-r_{i})/N.$
\end{lem}

Since we want to output a multiplicative approximation, we need $\Var(Y)$
to be small compared to $\E(Y)^{2}$, which means we need $r_{i}$
to be bounded from below. The choice of $K_{i}$ is precisely to ensure
this.
\begin{lem}
\label{lem:r_bound}For all $i=0,\ldots,m-1$, we have $\vol(K_{i+1})\le2\vol(K_{i})$. 
\end{lem}

\begin{proof}
We have 
\[
K_{i+1}=2^{(i+1)/n}B_{n}\cap K\subseteq2^{1/n}(2^{i/n}B_{n}\cap K)=2^{1/n}K_{i}.
\]
The lemma follows.
\end{proof}
Using the above lemmas in each phase, and keeping the total error
over $m$ phases smaller than a target $\varepsilon$ leads to a bound
of $N=O(m^{3}/\varepsilon^{2})$ samples per phase. The next lemma
which applies to any product of independent random variables gives
us the ``right'' bound.
\begin{lem}
\label{lem:prod_var}Let $Y=\prod_{i=0}^{m-1}Y_{i}$ where $Y_{i}\ge0$
are independent real-valued random variables. Then, 
\[
\frac{\Var(Y)}{\E(Y)^{2}}=\prod_{i=0}^{m-1}\left(1+\frac{\Var(Y_{i})}{\E(Y_{i})^{2}}\right)-1.
\]
\end{lem}

We can now prove the main theorem about the sample complexity of the
DFK volume algorithm.
\begin{thm}
For any input convex body $K$ satisfying $B_{n}\subseteq K\subseteq RB_{n}$
given by a membership oracle and any $\varepsilon>0$, with $N=O(\varepsilon^{-2}n\log R)$
samples per phase, with probability at least $3/4$, the DFK volume
algorithm outputs a multiplicative $(1+\varepsilon)$ approximation
to the volume of $K$. The total number of samples is thus $O^{*}(m^{2})=O^{*}(n^{2})$. 
\end{thm}

\begin{proof}
In each phase we have 
\[
\frac{\Var(Y_{i})}{\E(Y_{i})^{2}}=\frac{r_{i}(1-r_{i})}{Nr_{i}^{2}}=\frac{1-r_{i}}{Nr_{i}}\le\frac{1}{N}
\]
where we used the fact that $r_{i}\ge1/2$ from Lemma \ref{lem:r_bound}.
Now using Lemma \ref{lem:prod_var}, we have 
\[
\frac{\Var(Y)}{\E(Y)^{2}}\le\left(1+\frac{1}{N}\right)^{m}-1\le\frac{\varepsilon^{2}}{4}
\]
where we used that $(1+\frac{1}{N})^{m}-1\leq\frac{2m}{N}$ for $N\ge8m/\varepsilon^{2}$.
The theorem follows using Chebyshev's inequality.
\end{proof}
In \cite{kannan1997random}, it was shown that the amortized complexity
of uniform sampling is $O^{*}(n^{3})$ per sample when $\Omega(n^{2})$
samples are needed, giving an overall oracle complexity of $O^{*}(n^{5})$
for the entire algorithm. To make this rigorous, besides establishing
the amortized complexity of uniform sampling, we have to deal with
some nontrivial technical issues: (1) the distribution from which
samples are produced is close to the target uniform distribution but
not exactly uniform (2) samples are nearly independent but not perfectly
independent. Another important problem that will arise in establishing
the complexity of (approximate) uniform sampling is ``rounding''
which allows us to reduce the dependence of the sampling complexity
from polynomial in $R$ to polynomial in $\log R$. 

The quadratic sample complexity might appear to be nearly optimal
-- the final ratio being estimated could be as large as $R^{n}$
or $n^{n}$, which means that any sequence that uses at most $m$
phases estimates a ratio as large as $n^{n/m}$ needing as many samples,
and this means that $m=\Omega(n)$ is needed to get polynomial complexity
even. We will soon see that this intuitive reasoning is not valid. 

\section{Simulated Annealing}

In this section we introduce a sampling-based approach for optimization
and integration in high dimension. The main idea is to sample a sequence
of logconcave distributions, starting with one that is easy to integrate
and ending with the function whose integral is desired. This process
is known as \emph{simulated annealing. }The same high-level algorithm
can be used for optimization, volume computation/integration or rounding.
For integration, the desired integral of a function $f:\R^{n}\rightarrow\R_{+}$
can be expressed as the following telescoping product:

\[
\int_{\R^{n}}f=\int f_{0}\frac{\int f_{1}}{\int f_{0}}\frac{\int f_{2}}{\int f_{3}}\cdots\frac{\int f_{m}}{\int f_{m-1}}
\]
where $f_{0}$ is chosen to be a function with an easy integral (typically
a formula) and $f_{m}=f$. To compute any the ratio $\int g/\int h$
we can do the following: 
\begin{enumerate}
\item Sample $X$ from the distribution with density proportional to $h$. 
\item Use the estimator $Y=\frac{g(X)}{h(X)}$. 
\end{enumerate}
It is easy to establish the following fact about this estimator. 
\begin{lem}
For integrable functions $g,h$ with with support of $g$ contained
in the support of $h$ and $g(X)/h(X)$ defined everywhere in the
support of $h$, for $X$ drawn from the distribution with density
proportional to $h$ and $Y=g(X)/h(X)$, we have 
\[
\E(Y)=\frac{\int g}{\int h}.
\]
 
\end{lem}

Of course, the accuracy of the estimator depends on its variance.
Our goal is to choose a sequence of functions that are easy to sample
and lead to a small variance estimator. What sequence of functions
should we use?

In fact, we can view annealing as a general framework for efficient
statistical estimation by following a sequence of distributions, starting
with a particularly simple one and ending with the distribution of
interest. We state this ``meta''-algorithm below and later we will
specialize it to optimization and integration. 

\begin{algorithm2e}[H]

\caption{$\mathtt{SimulatedAnnealing}$}

\SetAlgoLined
\begin{enumerate}
\item For $i=0,\ldots,m$, define
\[
a_{i}=b\left(1+\frac{1}{\sqrt{n}}\right)^{i}\mbox{ and }f_{i}(x)=f(x)^{a_{i}}.
\]
\item Let $X_{0}^{1},\ldots,X_{0}^{k}$ be independent random points with
density proportional to $f_{0}$.
\item For $i=0,\ldots,m-1$: starting with $X_{i}^{1},\ldots,X_{i}^{k}$,
generate random points $X_{i+1}=\{X_{i+1}^{1},\ldots,X_{i+1}^{k}\}$;
update a running estimate $g$ based on these samples; update the
isotropic transformation using the samples.
\item Output the final $g$.
\end{enumerate}
\end{algorithm2e}

For optimization, the function $f_{m}$ is set to be a sufficiently
high power of $f$, the function to be maximized while $g$ is simply
the maximum objective value so far. For integration and rounding,
$f_{m}=f$, the target function to be integrated or rounded. For integration,
the function $g$ starts out as the integral of $f_{0}$ and is multiplied
by the ratio $\int f_{i+1}/\int f_{i}$ in each step. For rounding,
$g$ is simply the estimate of the isotropic transformation for the
current function.

For sampling and optimization, it is natural to ask why one uses a
sequence of functions rather than jumping straight to the target function
$f_{m}$. For optimizing a linear function $c^{T}x$ over a convex
set $K$, all we need to do is sample according to the density proportional
to $e^{-a(c^{T}x)}$ for a sufficiently large coefficient $a$ (recall
Theorem \ref{thm:OPTviaSampling}). The reason for using a sequence
is that the function $f_{m}$ and corresponding density $p_{m}$ can
be very far from the starting function or distribution, and hence
the mixing time can be high. The sequence ensures that samples from
the current function provide a ``warm'' start for sampling from
the next function. This is captured in the next lemma.
\begin{lem}
\label{thm:L2_warm_start} Let $p_{i}(x)=f_{i}(x)/\int f_{i}$ with
$f_{i}$ defined as in the annealing algorithm for some logconcave
function $f:\R^{n}\rightarrow\R_{+}$. Then a random sample $X\sim p_{i}$
satisfies
\[
\E_{p_{i}}\left(\frac{p_{i}(x)}{p_{i+1}(x)}\right)=O(1).
\]
\end{lem}

The underlying mathematical property behind this lemma is the following.
\begin{lem}
\label{lem:ratio_var} Let $f$ be a logconcave function in $\R^{n}.$
Then $Z(a)=a^{n}\int_{\R^{n}}f(x)^{a}$ is logconcave for $a\ge0$.
If $f$ has support $K$, then $Z(a)=a^{n}\int_{K}f(ax)$ is logconcave
for $a>0$.
\end{lem}

\begin{proof}
Consider the function 
\[
F(x,t)=f\left(\frac{x}{t}\right)^{t}
\]
for $x\in\R^{n}$ and $t>0$. Then $F$ is logconcave as verified
below. Consider $(x,t),(y,s)$ for $t,s>0$ and proceed as follows
using logconcavity: 
\begin{align*}
F\left(\frac{x+y}{2},\frac{t+s}{2}\right) & =f\left(\frac{x+y}{t+s}\right)^{\frac{t+s}{2}}\\
 & =f\left(\frac{x}{t}\cdot\frac{t}{t+s}+\frac{y}{s}\cdot\frac{s}{t+s}\right)^{\frac{t+s}{2}}\\
 & \ge f\left(\frac{x}{t}\right)^{\frac{t}{2}}f\left(\frac{y}{s}\right)^{\frac{s}{2}}=\sqrt{F(x,t)F(y,s)}.
\end{align*}

Now by the Prekopa-Leindler Theorem (see\ref{lem:marginal}), the
marginal of $F$ in $t$ is also logconcave. This marginal is 
\[
\int_{\R^{n}}F(x,t)\,dx=\int_{\R^{n}}f\left(\frac{x}{t}\right)^{t}\,dx=t^{n}\int_{\R^{n}}f(y)^{t}\,dy
\]
where we made the substitution $y=x/t$.
\end{proof}
Note that in the simulated annealing algorithm we choose $\frac{a_{\{i+1\}}}{a_{i}}=(1+\frac{1}{\sqrt{{n\}}}})$
since this guarantees that the distribution with density proportional
to $f_{i}$ has $\chi^{2}$-distance $O(1)$ to the distribution with
density proportional to $f_{i+1}$. This will be useful later in bounding
$\text{{var}(Y)/\ensuremath{\mathbb{E}[Y]^{2}}}$ of the random variable
given by $Y=f_{i+1}(X)/f_{i}(X)$ for a random sample $X$ drawn from
the distribution with density proportional to $f_{i}$.
